\section*{Capítulo \ref{ch:g8:equations} --- Cálculo literal e equações}

\begin{solution}{exo:g8:equations:1}
$5(2x + 3) = 10x + 15$.

$-3(4x - 1) = -12x + 3$.

$2(3x - 2) + 4(x + 5) = 6x - 4 + 4x + 20 = 10x + 16$.
\end{solution}

\begin{solution}{exo:g8:equations:2}
$(x+2)(x+7) = x^2 + 7x + 2x + 14 = x^2 + 9x + 14$.

$(3x+1)(2x-5) = 6x^2 - 15x + 2x - 5 = 6x^2 - 13x - 5$.

$(x-4)(x-6) = x^2 - 6x - 4x + 24 = x^2 - 10x + 24$.
\end{solution}

\begin{solution}{exo:g8:equations:3}
$4 \times 3 - 5 = 7$: sim. \quad
$2 \times 3 + 1 = 7$ e $3 + 4 = 7$: sim. \quad
$3^2 = 9$ e $6 \times 3 - 9 = 9$: sim --- $x = 3$ resolve as três.
\end{solution}

\begin{solution}{exo:g8:equations:4}
$x + 9 = 4$: $x = -5$ (confere: $-5 + 9 = 4$).

$5x = -35$: $x = -7$.

$\frac x3 = 8$: $x = 24$.

$2x - 7 = 13$: $2x = 20$, $x = 10$ (confere: $20 - 7 = 13$).
\end{solution}

\begin{solution}{exo:g8:equations:5}
$6x + 5 = 2x + 25$: $4x = 20$, então $x = 5$.

$9 - 4x = x - 11$: $20 = 5x$, então $x = 4$.

$3(x + 2) = 2x + 11$: $3x + 6 = 2x + 11$, então $x = 5$.
\end{solution}

\begin{solution}{exo:g8:equations:6}
$8x - 12 = 5x + 30$, então $3x = 42$ e $x = 14$. Confere:
$4(28 - 3) = 100$ e $5(14 + 6) = 100$.
\end{solution}

\begin{solution}{exo:g8:equations:7}
``$3x - 8 = 19$'': $3x = 27$, $x = 9$.

``$2x + 5 = x + 12$'': $x = 7$.
\end{solution}

\begin{solution}{exo:g8:equations:8}
Largura $x$, comprimento $x + 5$; \omterm{def:g6:measure:perimeter}{perímetro}
$2\bigl(x + (x + 5)\bigr) = 4x + 10 = 58$, então $4x = 48$ e $x = 12$:
largura $12$ cm, comprimento $17$ cm (confere:
$2 \times (12 + 17) = 58$).
\end{solution}

\begin{solution}{exo:g8:equations:9}
Os números são $n - 1$, $n$, $n + 1$, de \omterm{def:g1:addition:def}{soma} $3n = 141$, então
$n = 47$: os números são $46$, $47$, $48$.
\end{solution}

\begin{solution}{exo:g8:equations:10}
$3x + 4 > 19$: $3x > 15$, então $x > 5$.

$7 - 2x \geq 1$: $-2x \geq -6$, e dividir por $-2$ inverte:
$x \leq 3$.

$5(x-1) < 2x + 7$: $5x - 5 < 2x + 7$, então $3x < 12$ e $x < 4$.
\end{solution}

\begin{solution}{exo:g8:equations:11}
Depois de $n$ meses, Mia tem $140 + 15n$ e Tomás tem $200 + 10n$. Mia
fica na frente quando
\[
140 + 15n > 200 + 10n
\ \Longleftrightarrow\
5n > 60
\ \Longleftrightarrow\
n > 12 :
\]
a partir do $13^{\text{o}}$ mês.
\end{solution}

\begin{solution}{exo:g8:equations:12}
Multiplique os dois lados por $20$:
\[
5(x + 3) = 4(2x - 1)
\ \Longrightarrow\
5x + 15 = 8x - 4
\ \Longrightarrow\
19 = 3x
\ \Longrightarrow\
x = \frac{19}{3}.
\]
Confere: $\frac{19/3 + 3}{4} = \frac{28/3}{4} = \frac{28}{12} =
\frac73$, e $\frac{2 \times 19/3 - 1}{5} = \frac{35/3}{5} =
\frac{35}{15} = \frac73$. Iguais: a solução é $\frac{19}{3}$.
\end{solution}

\begin{solution}{pb:g8:equations:1}
\textbf{1.} Distributividade dupla (\cref{thm:g8:equations:expand})
com $c = a$ e $d = b$:
\[
(a + b)(a + b) = a^2 + ab + ba + b^2 = a^2 + 2ab + b^2 ,
\]
já que $ab$ e $ba$ são o mesmo \omterm{def:g2:mult:def}{produto}.

\textbf{2.} O mesmo desenvolvimento, com atenção aos sinais:
\[
(a - b)(a - b) = a^2 - ab - ba + b^2 = a^2 - 2ab + b^2 ,
\]
\[
(a + b)(a - b) = a^2 - ab + ba - b^2 = a^2 - b^2 :
\]
desta vez os termos cruzados $-ab$ e $+ba$ se cancelam em vez de
dobrar.

\textbf{3.} O quadrado de lado $a + b$ se corta em quatro peças: um
quadrado $a \times a$ (o $a^2$), um quadrado $b \times b$ (o $b^2$) e
\emph{dois} retângulos $a \times b$ (juntos, o $2ab$). A identidade
$(a+b)^2 = a^2 + 2ab + b^2$ é a afirmação de que as quatro \omterm{def:g6:measure:area}{áreas}
preenchem o quadrado grande --- é o desenho da distributividade dupla
com os dois lados divididos do mesmo jeito.

\textbf{4.} Com os \omterm{def:g2:mult:def}{produtos} notáveis:
\begin{align*}
101^2 &= (100 + 1)^2 = 10\,000 + 2 \times 100 + 1 = 10\,201, \\
99^2 &= (100 - 1)^2 = 10\,000 - 200 + 1 = 9\,801, \\
49 \times 51 &= (50 - 1)(50 + 1) = 50^2 - 1^2 = 2\,499 .
\end{align*}

\textbf{5.} Para $a = 3$ e $b = 4$: $(3 + 4)^2 = 49$, mas
$3^2 + 4^2 = 9 + 16 = 25$. A fórmula de Zoe esquece os dois retângulos
$a \times b$ do desenho --- os $2ab = 24$ que faltam, e de fato
$25 + 24 = 49$.

\textbf{6.} $1$; $1 + 3 = 4$; $1 + 3 + 5 = 9$; $1 + 3 + 5 + 7 = 16$;
$1 + 3 + 5 + 7 + 9 = 25$. Os resultados são os quadrados perfeitos
$1^2, 2^2, 3^2, 4^2, 5^2$.

\textbf{7.} Os \omterm{def:g3:numbers:evenodd}{números ímpares} começam em $1$ e sobem de $2$ em $2$: o
$n$-ésimo é alcançado depois de $n - 1$ passos de $2$ a partir de $1$,
então vale $1 + 2(n - 1) = 2n - 1$. Confere: $n = 1, 2, 3$ dão
$1, 3, 5$.

\textbf{8.} Pelo primeiro \omterm{def:g2:mult:def}{produto} notável,
\[
(n + 1)^2 - n^2 = n^2 + 2n + 1 - n^2 = 2n + 1 :
\]
a \omterm{ex:g1:subtraction:difference}{diferença} entre dois quadrados consecutivos é exatamente o
$(n+1)$-ésimo \omterm{def:g3:numbers:evenodd}{número ímpar} (questão 7 com $n + 1$ no lugar de $n$:
$2(n+1) - 1 = 2n + 1$).

\textbf{9.} Construa os quadrados um depois do outro. Comece com uma
única casinha: $1 = 1^2$. Para passar do quadrado de lado $1$ ao de
lado $2$, acrescente $2 \times 1 + 1 = 3$ casinhas (questão 8):
$1 + 3 = 2^2$. Para passar ao de lado $3$, acrescente
$2 \times 2 + 1 = 5$ casinhas: $1 + 3 + 5 = 3^2$. Cada \omterm{def:g3:numbers:evenodd}{número ímpar}
novo é exatamente a camada em L que completa o próximo quadrado, então,
depois de somar os $n$ primeiros \omterm{def:g3:numbers:evenodd}{ímpares}, o quadrado de lado $n$ está
completo: $1 + 3 + \dots + (2n - 1) = n^2$.

\textbf{10.} Os \omterm{def:g3:numbers:evenodd}{números ímpares} de $1$ a $99$ são os $50$ primeiros
($2n - 1 = 99$ dá $n = 50$), então
\[
1 + 3 + 5 + \dots + 99 = 50^2 = 2\,500 .
\]
E $1 + 3 + \dots + (2n - 1) = 400 = 20^2$ exige $n = 20$: os
\emph{vinte} primeiros \omterm{def:g3:numbers:evenodd}{números ímpares} ($1$ até $39$) somam $400$.

\textbf{11.} $x^2 + 6x + 9 = x^2 + 2 \times 3 \times x + 3^2 =
(x + 3)^2$: o primeiro \omterm{def:g2:mult:def}{produto} notável com $a = x$ e $b = 3$. Então a
\omterm{def:g8:equations:def}{equação} fica $(x + 3)^2 = 0$. Um quadrado só é \omterm{def:g1:counting:zero}{zero} quando o próprio
número é \omterm{def:g1:counting:zero}{zero} (\cref{thm:g8:negprod:rules}: um número não nulo tem
distância não nula, então o quadrado dele também tem distância não
nula). Logo $x + 3 = 0$: a única solução é $x = -3$.

\textbf{12.} Pelo \cref{thm:g8:negprod:rules}, a distância de um
\omterm{def:g2:mult:def}{produto} é o \omterm{def:g2:mult:def}{produto} das distâncias. Se os dois fatores forem não nulos,
as duas distâncias são não nulas, então a distância do \omterm{def:g2:mult:def}{produto} é não
nula: o \omterm{def:g2:mult:def}{produto} não pode ser $0$. Um \omterm{def:g2:mult:def}{produto} nulo, portanto, obriga
pelo menos um fator nulo. Para $(x - 2)(x + 5) = 0$: ou $x - 2 = 0$, ou
$x + 5 = 0$, o que dá as duas soluções $x = 2$ e $x = -5$.

\textbf{13.} $x^2 = 49$ vira $x^2 - 49 = 0$, e o terceiro \omterm{def:g2:mult:def}{produto}
notável o fatora:
\[
x^2 - 49 = x^2 - 7^2 = (x - 7)(x + 7) = 0 .
\]
Pela questão 12, $x = 7$ ou $x = -7$: \emph{duas} soluções. Toda
\omterm{def:g8:equations:def}{equação} resolvida até aqui neste capítulo era do primeiro grau e tinha
exatamente uma solução; um quadrado na \omterm{def:g8:equations:def}{equação} pode produzir duas.

\textbf{14.} Passe tudo para o lado esquerdo: $x^2 - 6x + 9 = 0$, e
reconheça o segundo \omterm{def:g2:mult:def}{produto} notável:
$x^2 - 2 \times 3 \times x + 3^2 = (x - 3)^2$. Então a \omterm{def:g8:equations:def}{equação} é
$(x - 3)^2 = 0$, o que obriga $x - 3 = 0$ (como na questão 11):
$x = 3$ é a única solução --- a conferência do
\cref{exo:g8:equations:3} a encontrou, e a fatoração demonstra que não
há outras.

\textbf{15.} Pelo terceiro \omterm{def:g2:mult:def}{produto} notável,
\[
2026^2 - 2025^2 = (2026 - 2025)(2026 + 2025)
= 1 \times 4051 = 4\,051 .
\]
Em geral, $(n+1)^2 - n^2 = (n + 1 - n)(n + 1 + n) = 2n + 1$: a
\omterm{ex:g1:subtraction:difference}{diferença} de dois quadrados consecutivos é a \omterm{def:g1:addition:def}{soma} dos dois números ---
o mesmo $2n + 1$ da questão 8, obtido agora pelo terceiro \omterm{def:g2:mult:def}{produto}
notável em vez do primeiro.
\end{solution}
